\chapter{Cronograma de execução}\label{cap:cronograma}


O cronograma proposto contempla as etapas necessárias para o desenvolvimento e validação da metodologia descrita, distribuídas ao longo de seis meses. As atividades foram planejadas de modo sequencial, com sobreposição controlada entre as fases de processamento, análise e integração, visando otimizar o tempo de execução e garantir consistência nos resultados.

\begin{table}[H]
\centering
\caption{Cronograma das atividades do projeto}
\begin{tabular}{|p{5cm}|c|c|c|c|c|c|}
\hline
\textbf{Atividade} & \textbf{Mês 1} & \textbf{Mês 2} & \textbf{Mês 3} & \textbf{Mês 4} & \textbf{Mês 5} & \textbf{Mês 6} \\
\hline
Levantamento da base de dados & X & X & & & & \\
\hline
Anonimização dos textos & & X & X & & & \\
\hline
Segmentação semântica dos textos  & & X & X & & & \\
\hline
Extração de entidades clínicas (\textit{Entities Extraction}) & & & X & X & & \\
\hline
Ligação de relacionamentos e construção do grafo  & & & & X & X & \\
\hline
Consultas e recuperação de informação & & & & X & X & \\
\hline
Avaliação dos resultados e ajustes metodológicos & & & & & X & X \\
\hline
Redação do relatório técnico e conclusões & & & & & & X \\
\hline
\end{tabular}
\end{table}

As atividades serão executadas de forma incremental, permitindo revisões contínuas conforme os resultados forem sendo obtidos. O cronograma foi estruturado para garantir a integração gradual entre as etapas de pré-processamento, extração semântica e representação em grafo, assegurando a qualidade e a coerência das informações processadas.
